\chapter{Discrete choice models}

\section{Supported models}
\label{sec:models}

\amlet{} allows the estimation of multinomial (or conditional) logit
models as well as mixed logit models.
When use as a standalone program \amlet{} automatically determines the model type from the utility description.
The executable version of \amlet{} currently supports linear-in-parameters utilities only, i.e., for $i = 1,\ldots,I$, $j = 1,\ldots,J$,
\[
U_{ij}(x_{ij}) = \sum_{k = 1}^{n} \beta_{ik}x_{ijk} = \beta^T x_{ij},
\]
where $n$ is the number of dependant variables. Each component of the
vector $\beta$ can be fixed, constant or random, as summarized in
Table~\ref{tab:parameters}.

\begin{table}[htb]
\begin{center}
\begin{tabular}{|c|c|}
\hline
\verb|FIXED| & fixed coefficient \\
\hline
\verb|CONSTANT| & constant (non-random) coefficient \\
\hline
\verb|LOGNORMAL| & lognormally-distributed coefficient \\
\hline
\verb|NORMAL| & normally-distributed coefficient \\
\hline
\verb|SPLINE| & spline-distributed coefficient \\
\hline
\verb|TRIANGULAR| & triangularly-distributed coefficient \\
\hline
\verb|UNIFORM| & uniformly-distributed coefficient \\
\hline
\end{tabular}
\caption{parameters types}
\label{tab:parameters}
\end{center}
\end{table}

\section{Supported distributions}

The following distributions are currently supported for the
coefficient of utilities in a mixed logit model and for observations
when generating synthetic data.
\begin{description}
\item[Lognormal]
lognormally distributed coefficient; the coefficient is calculated as
\[e^{\mu + \sigma X}\] where $X \sim N(0,1)$, and $\mu$ and $\sigma$
are the parameters to estimate. Therefore $\mu$ and $\sigma$ are
respectively the mean and the standard deviation of the logarithm of
the coefficient.

%A lognormal distribution with parameters $\mu$ and $\sigma$ has median
%$e^\mu$, mean $e^{\mu+\sigma^2/2}$, and standard deviation $e^{\sigma^2-1}$.

Since lognormally distributed coefficients are always positive; you
have to enter the negative of the explanatory variable in your
dataset if necessary.
\item[Normal] normally distributed coefficient, with the mean and
standard deviation being estimated.
\item[Spline] \textit{To check.}
\item[Triangular]
Triangularly distributed coefficient, with the mean and ``spread'' being
estimated. A triangular distribution with mean $b$ and spread $a$ has
a zero density outside $[b-a, b+a]$, while it increases linearly from
$b-a$ to $b$ and it decreases linearly from $b$ to $b+a$.
\item[Uniform]
uniformly distributed coefficient, with the mean $b$ and "spread" $a$
being estimated. A uniformly distributed coefficient of mean $b$ and
spread $a$ has therefore the distribution of $U[b-a,b+a]$.
\end{description}

You can also use the previous distributions to generate synthetic
data, as well as the following one:
\begin{description}
\item[Bernouilli] an observation can be drawn from a Bernouilli of
parameter $p$ ($0 \leq p \leq 1.0$), that has to be specified.
\end{description}

The used distributions are specified in the variables
\verb|btype| and \verb|bgentype|
as set of identifiers.
The truncated normal distributions (\verb|truncated_normal| and
\verb|truncated_normal_b|) and the Bernouilli distribution can only be
used in \verb|bgentype|, to generate synthetic data.

%\section{Alternatives sampling}

\section{Sampling methods}
\label{sec:sampling}

Multi-dimensional integrals are approximated by drawings points from
the underlying distributions of the random parameters.
Standard Monte-Carlo draws can be used as well as some quasi-Monte
Carlo techniques:
\begin{itemize}
\item
Halton sequences~\cite{Halt60},
\item
Lattice nets~\cite{MungLEcuBastCiriTuff12},
\item
Modified Latin Hypercube Sampling as proposed in Hess et al.~\cite{HessTraiPola04},
\item
Sobol sequences, following the recommendations of Joe and Kuo~\cite{JoeKuo03}.
\end{itemize}
These various quasi-Monte Carlo sequences can be generated by using
the library \oratio().
The generation method can be defined by with the variable \verb|sampling|, and is set by default to the standard Monte-Carlo approach.
Possible values are
summarized in Table~\ref{tab:samplingmethod}.

\begin{table}[htb]
\begin{center}
\begin{tabular}{|c|c|}
\hline
Pseudo-random numbers & \verb|monte_carlo| \\
\hline
Halton sequences & \verb|halton| \\
\hline
Modified Latin Hypercube Sampling & \verb|mlhs| \\
\hline
Rank-1 Lattice & \verb|lattice| \\
\hline
Sobol sequences & \verb|sobol| \\
\hline
\end{tabular}
\caption{sampling methods}
\label{tab:samplingmethod}
\end{center}
\end{table}

The number of random draws to use per individual is defined with the
variable \verb|nrep|, and is set by default to 1000. A higher value is
sometimes needed, especially if the accuracy and bias on the objective
are important. To use more draws, say 10000 per individual, write
\begin{verbatim}
nrep = 10000;
\end{verbatim}
When lattices are used, \amlet{} uses the tabulated rank-1 lattice with
the largest number of points (a prime number) not greater than
\verb|nrep|, and sets \verb|nrep| to this number, for instance 887 for
\verb|nrep = 1000|.

The methods with a variable sample size (\verb|mcbtrda| and
\verb|mcrtrda|) estimate the simulation error as for independent
draws. With quasi-Monte Carlo draws, \amlet{} therefore replaces them by
the same methods with a fixed sample size (\verb|mcbtr| and
\verb|mcrtr|), with a warning.
While Halton are the most popular sequences in discrete choice theory, we recommend the use of Sobol sequences of rank-1 lattices.
However, the advantage of quasi-Monte Carlo sequences is not clear, as shown in \cite{MungLEcuBastCiriTuff12}.

\subsection{Halton sequences}

A popular alternative to pseudo-random draws in the field of mixed
logit modelling is the use of Halton sequences
(Halton~\cite{Halt60}). While we strongly recommend the use of Sobol
sequences instead of Halton ones, it is possible to select this
approach by correspondly setting the variable \verb|sampling|:
\begin{verbatim}
sampling = halton;
\end{verbatim}
One major drawback of Halton sequences is the correlation between
sequences associated to successive dimensions when the dimensionality
is quite high (order of 10 or superior). Two techniques have been
proposed to break such correlations: scrambling (Bhat~\cite{Bhat03})
and shuffling (Hess et al.~\cite{HessPolaDaly03}). Both of them can be used in
\amlet{}. It is also possible to randomize the sequences by randomly
shifting the sequences, for each dimension, as proposed by
Bhat~\cite{Bhat03}. While all these techniques are not strictly
speaking randomizations methods, all of them can be set using the
variable \verb|randomization|, whose default value is \verb|none|, as
summarized in Table~\ref{tab:halton_randomization}.

\begin{table}[htb]
\begin{center}
\begin{tabular}{|c|c|}
\hline
\verb|none| & no randomization \\
\hline
\verb|scrambling| & scramble the Halton sequences \\
\hline
\verb|shifting| & randomly shift the Halton sequences \\
\hline
\verb|shifting_scrambling| & randomly shift and scramble the Halton
sequences \\
\hline
\verb|shuffling| & randomly shuffle the Halton sequences \\
\hline
\verb|shifting_shuffling| & randomly shift and shuffle the Halton sequences \\
\hline
\end{tabular}
\caption{Randomization strategies for Halton sequences}
\label{tab:halton_randomization}
\end{center}
\end{table}

In \amlet{}, all individuals use the same Halton points, each with its
own random shift, so that \verb|none| and \verb|shifting| give the same
draws. \verb|scrambling| applies the Braaten-Weller permutations (up to
dimension 16), and \verb|shuffling| shuffles the order of the points of
each individual, dimension by dimension. Up to the 2026 rebuild,
\amlet{} ignored the variable \verb|randomization|; it still does with
\verb|legacy_rng = true|.

\subsection{Sobol sequences}

\subsection{Rank-1 lattices}

\section{Pseudo-random numbers generators}
\label{sec:rng}

Each method requires the use of a uniform random numbers generator,
whose quality is important to ensure quality of the results.
Even if a quasi-Monte Carlo is used, the randomization part also involves the
use of a pseudo-random numbers generator.
The choice of a good one is still a tricky question (see for instance the
relevant discussion in L'Ecuyer~\cite{LEcu97} and McCullough~\cite{McCu06}).
Unlike many econometric packages, as \limdep, \amlet{} allows the user
to choose between several random number generators, defined in the
random numbers generator library.
Three generators are currently available: the standard minimal one,
proposed by Park and Miller~\cite{ParkMill88}, the MRG32K3a generator, due to L'Ecuyer, and the generator MT19937 designed by
Matsumoto and Nishimura~\cite{MatsNish98}.
To date no significant differences have been observed when using one
technique instead of another one, while it is known that the sandard minimal generator has severe statistical deficiencies.
By default we use the MRG32K3a generator that benefits from strong theoretical properties and is well tested, and our implementation allows to consider substreams.
The choice of the random number generator can be defined with the
variable \verb|rangen|. Possible values are summarized in
Table~\ref{tab:rangen}.

\begin{table}[htb]
\begin{center}
\begin{tabular}{|c|c|}
\hline
Standard minimal & \verb|standard| \\
\hline
L'Ecuyer's generator & \verb|ecuyer| \\
\hline
MRG32K3a (default) & \verb|mrg32k3a| \\
\hline
MT19937 & \verb|mt19937| \\
\hline
\end{tabular}
\caption{random number generators}
\label{tab:rangen}
\end{center}
\end{table}

Up to the 2026 rebuild of \amlet{} 0.13.1, the MRG32K3a generator of
\oratio{} 0.7.5 differed from L'Ecuyer's: its first component used
$x_{n-1}$ instead of $x_{n-2}$, and the jumps to the next stream or
substream were not reduced modulo $m$. It now produces L'Ecuyer's
MRG32K3a sequences, streams and substreams, so that simulated results
differ from those of \amlet{} 0.13.1. To reproduce them, set
\begin{verbatim}
legacy_rng = true;
\end{verbatim}
in the driver file. The legacy generator is not MRG32K3a and its
statistical properties are unknown: use it only to reproduce earlier
results.

Each random generator produces a sequence of points that mimics draws
from an uniform distribution on $(0,1)$.. These sequences are however
deterministic due to their numerical nature, so they can be
reproduced when some initial conditions are fixed. Typically, they can
be expressed by a single strictly positive integer, known as the {\em
  seed}. This value can be defined in the driver file by means of the
variable \verb|seed1|, for example:
\begin{verbatim}
seed1 = 12345;
\end{verbatim}
Most of the time, the seed can be recovered at any point during the
random draws generation process, allowing the replication of a
sequence starting at any draw of the initial sequence. This is useful
for instance when dealing with a newly generated synthetic population, where the same random generator is used for creating the population and the random draws.
In this
case, the output file will contain both the initial seed and the seed
after the population generation. Note however that the MT19937 cannot
be restarted with a single seed, except the initial one.

The default value for \verb|seed1| is 12345.
If a seed value of 0 is used, \amlet{} will determine a seed based on the computer clock, that will be used to initiate the pseudo-random draws sequence.

For more information, we refer the reader to the documentation
accompying the \oratio{} library.

\section{Results validation}

The source file \verb|validation.c|, associated to the header file
\verb|validation.h|, contains basic validation routines concerned with
validation of the estimation results, on a statistical point of view.
The validation is currently quite simple, as it simply repeats
\verb|nsim| times the simulation of the objective function, using
a different set of random realizations for each simulation.
The mean, standard deviation, and a confidence interval at 0.9 for the
objective function are given on output.
In order to activate the validation function, you need to set the
variable \verb|validation| to \verb|true| in the driver file:
\begin{verbatim}
validation = true;
\end{verbatim}
